\documentclass[11pt]{article}
\usepackage{amsmath, amssymb}
\usepackage[margin=1in]{geometry}
\usepackage{booktabs}

\title{Upper-Level Fairness Reformulations for First-Iteration Feasibility}
\author{FairLoadDelivery Package Documentation}
\date{March 2026}

\begin{document}
\maketitle

\section{Background}

The bilevel Fair Load Delivery Problem (FLDP) iterates between a lower-level Minimum Load Delivery (MLD) problem and an upper-level fairness optimization.
At each iteration~$k$, the lower level returns load-shed values $\mathbf{p}^{\mathrm{shed}}$ and a Jacobian $\mathbf{J} = \partial \mathbf{p}^{\mathrm{shed}} / \partial \mathbf{w}$ via implicit differentiation.
The upper level then updates the weight vector $\mathbf{w}$ to improve a chosen fairness metric, using a first-order Taylor approximation of the lower-level response:
\begin{equation}\label{eq:taylor}
    \hat{p}_i^{\mathrm{shed}}(\mathbf{w})
    = p_i^{\mathrm{shed,(prev)}}
      + \sum_{j=1}^{n} J_{ij} \bigl(w_j - w_j^{\mathrm{(prev)}}\bigr),
    \qquad i = 1,\dots,n.
\end{equation}

All upper-level formulations share the following weight constraints:
\begin{align}
    1 &\le w_j \le 10, & j &= 1,\dots,n, \label{eq:wbounds}\\
    |w_j - w_j^{\mathrm{(prev)}}| &\le \Delta, & j &= 1,\dots,n, \label{eq:trust}
\end{align}
where $\Delta = 0.5$ is the trust-region radius.
The trust region is two-sided, limiting both increases and decreases to~$\Delta$ per iteration.
This ensures the first-order Taylor approximation~\eqref{eq:taylor} remains accurate, since the Jacobian is only valid locally.

\paragraph{Correction from one-sided trust region.}
The original implementation used a one-sided constraint ($w_j - w_j^{\mathrm{(prev)}} \le \Delta$), which only limited weight \emph{increases}.
On the first iteration with $\mathbf{w}^{\mathrm{(prev)}} = 10 \cdot \mathbf{1}$, weights could decrease by up to~$9$ units (to the lower bound of~$1$), far exceeding the region where the Taylor approximation is valid.
This caused fairness objectives (Palma, equality-min, min-max) to produce degenerate solutions by uniformly decreasing all weights, inflating predicted load shed to~$100\%$ across all loads.

\section{Proportional Fairness}\label{sec:proportional}

\subsection{Original Formulation}

Proportional fairness (Nash bargaining, $\alpha$-fairness with $\alpha=1$) maximizes the sum of logarithms of load served.
The original formulation introduced auxiliary variables $s_i$ for load served:
\begin{align}
    \max_{\mathbf{w},\,\mathbf{s}} \quad & \sum_{i=1}^{n} \log(s_i) \\
    \text{s.t.} \quad
    & s_i = p_i^{\mathrm{ref}} - \hat{p}_i^{\mathrm{shed}}(\mathbf{w}), & i &= 1,\dots,n, \label{eq:prop-eq}\\
    & s_i \ge \varepsilon, & i &= 1,\dots,n, \label{eq:prop-lb}\\
    & \eqref{eq:wbounds},\;\eqref{eq:trust}, \notag
\end{align}
where $p_i^{\mathrm{ref}}$ is the reference (pre-shedding) demand for load~$i$ and $\varepsilon = 10^{-6}$.

\subsection{Issues}

\begin{enumerate}
    \item \textbf{Infeasibility from equality + lower bound.}
    Constraint~\eqref{eq:prop-eq} forces $s_i = p_i^{\mathrm{ref}} - \hat{p}_i^{\mathrm{shed}}$.
    If any load is fully shed on the first iteration ($p_i^{\mathrm{shed,(prev)}} \approx p_i^{\mathrm{ref}}$), the Taylor approximation~\eqref{eq:taylor} may not be able to reduce $\hat{p}_i^{\mathrm{shed}}$ enough within the trust region to satisfy $s_i \ge \varepsilon$.

    \item \textbf{Index mismatch between $\mathbf{p}^{\mathrm{ref}}$ and $\mathbf{p}^{\mathrm{shed}}$.}
    The original implementation built $\mathbf{p}^{\mathrm{ref}}$ from all loads in the network data dictionary (sorted by ID), while $\mathbf{p}^{\mathrm{shed}}$ from the lower level may exclude certain loads (e.g., source bus loads).
    This caused $p_i^{\mathrm{ref}}$ to correspond to a different load than $p_i^{\mathrm{shed}}$, leading to $p_i^{\mathrm{ref}} - p_i^{\mathrm{shed}} < 0$ and \texttt{log(negative)} evaluations, which Ipopt reports as \texttt{INVALID\_MODEL}.
\end{enumerate}

\subsection{Reformulation: Shifted Logarithm with Aligned Indexing}

We make two changes:
\begin{enumerate}
    \item \textbf{Signature change:} The function now receives a pre-aligned demand vector $\mathbf{p}^{\mathrm{ref}}$ (built by the caller from \texttt{pshed\_ids}), rather than constructing it internally from the math dictionary. This matches the pattern used by the Palma ratio function.
    \item \textbf{Shifted log with variable-based formulation:} We use $\hat{p}_i^{\mathrm{shed}}$ as a JuMP variable (linked to the Taylor expression by equality constraints) and apply the shifted logarithm:
\end{enumerate}
\begin{align}
    \max_{\mathbf{w},\,\hat{\mathbf{p}}^{\mathrm{shed}}} \quad & \sum_{i=1}^{n} \log\!\bigl(p_i^{\mathrm{ref}} - \hat{p}_i^{\mathrm{shed}} + c\bigr) \\
    \text{s.t.} \quad
    & \hat{p}_i^{\mathrm{shed}} = p_i^{\mathrm{shed,(prev)}} + \sum_{j=1}^{n} J_{ij} \Delta w_j, & i &= 1,\dots,n, \\
    & \eqref{eq:wbounds},\;\eqref{eq:trust}, \notag
\end{align}
where $c = 10^{-6}$ is a uniform shift applied to all loads' served values.
The shift ensures $\log(\cdot)$ remains defined even when a load is fully shed ($\hat{p}_i^{\mathrm{shed}} = p_i^{\mathrm{ref}}$ gives $\log(c)$, which is finite).
Ipopt's \texttt{warm\_start\_init\_point} option is enabled with starting values $\mathbf{w} = \mathbf{w}^{\mathrm{(prev)}}$ and $\hat{\mathbf{p}}^{\mathrm{shed}} = \mathbf{p}^{\mathrm{shed,(prev)}}$, ensuring the initial log evaluation is valid.

\section{Min-Max Load Shed}\label{sec:minmax}

\subsection{Original Formulation}

The min-max objective minimizes the worst-case (maximum) load shed across all loads using an auxiliary variable~$t$:
\begin{align}
    \min_{\mathbf{w},\,t} \quad & t \\
    \text{s.t.} \quad
    & t \ge \max_i \hat{p}_i^{\mathrm{shed}}(\mathbf{w}), \label{eq:mm-max}\\
    & t \ge 1, \notag \\
    & \eqref{eq:wbounds},\;\eqref{eq:trust}. \notag
\end{align}

\subsection{Issues}

Two problems made this formulation unreliable:
\begin{enumerate}
    \item The $\max(\cdot)$ in constraint~\eqref{eq:mm-max} was implemented using Julia's \texttt{maximum()} function on the JuMP expression array, which evaluates non-smoothly at model construction time rather than creating proper optimization constraints. This causes incorrect or infeasible behavior.
    \item The lower bound $t \ge 1$ is unnecessarily restrictive; when all load-shed values are small, the problem may be infeasible.
\end{enumerate}

\subsection{Reformulation: Epigraph}

We replace the non-smooth \texttt{maximum()} with individual constraints (the standard epigraph reformulation of the infinity norm) and relax $t \ge 1$ to $t \ge 0$:
\begin{align}
    \min_{\mathbf{w},\,t} \quad & t \\
    \text{s.t.} \quad
    & t \ge \hat{p}_i^{\mathrm{shed}}(\mathbf{w}), & i &= 1,\dots,n, \\
    & t \ge 0, \notag \\
    & \eqref{eq:wbounds},\;\eqref{eq:trust}. \notag
\end{align}
This is mathematically equivalent to $\min \|\hat{\mathbf{p}}^{\mathrm{shed}}\|_\infty$ but uses smooth, solver-compatible constraints.
Note: JuMP does not support \texttt{norm(x, Inf)} as a nonlinear operator for any solver, so the epigraph form is required.

\section{Palma Ratio}\label{sec:palma}

\subsection{Original Formulation}

The original Palma ratio formulation minimized the ratio of the top~10\% to bottom~40\% of sorted load-shed values using the Charnes-Cooper transformation:
\begin{align}
    \min_{\mathbf{w},\,\sigma} \quad & \sum_{i \in \mathrm{Top10}} \mathrm{sorted}_i \cdot \sigma \\
    \text{s.t.} \quad
    & \sum_{i \in \mathrm{Bot40}} \mathrm{sorted}_i \cdot \sigma = 1, \\
    & \sigma \ge \varepsilon, \notag \\
    & \text{permutation, McCormick, trust region constraints.} \notag
\end{align}

\subsection{Issues}

\begin{enumerate}
    \item \textbf{Wrong objective target.} Minimizing the Palma ratio pushes it below~$1$, meaning the top~10\% sheds \emph{less} than the bottom~40\%---reverse inequality, not fairness. The desired target is Palma ratio $= 1$ (proportional balance).
    \item \textbf{No efficiency incentive.} The ratio objective has no penalty for increasing total load shed. The optimizer can equalize shedding by inflating all values toward~$100\%$, which achieves ratio $\approx 1$ but is a degenerate solution.
\end{enumerate}

\subsection{Reformulation: Target Ratio = 1}

We replace the Charnes-Cooper transformation with a squared-deviation objective targeting Palma ratio $= 1$:
\begin{align}
    \min_{\mathbf{w}} \quad & \Bigl(\sum_{i \in \mathrm{Top10}} \mathrm{sorted}_i - \sum_{i \in \mathrm{Bot40}} \mathrm{sorted}_i\Bigr)^2 \\
    \text{s.t.} \quad & \text{permutation, McCormick, trust region constraints.} \notag
\end{align}
This eliminates the Charnes-Cooper variable $\sigma$ and its quadratic constraint, simplifying the formulation.
When the top~10\% and bottom~40\% sums are equal, the objective is zero and the Palma ratio is exactly~$1$.
The Charnes-Cooper code is retained (commented out) for cases where direct ratio minimization is desired.

\section{Equality-Min Fairness}\label{sec:equality}

\subsection{Original Formulation}

The equality-min objective seeks to equalize all load-shed values by minimizing a common target~$t$ subject to hard equality:
\begin{align}
    \min_{\mathbf{w},\,t} \quad & t \\
    \text{s.t.} \quad
    & t = \hat{p}_i^{\mathrm{shed}}(\mathbf{w}), & i &= 1,\dots,n, \label{eq:eq-hard}\\
    & t \ge 0, \notag \\
    & \eqref{eq:wbounds},\;\eqref{eq:trust}. \notag
\end{align}

\subsection{Infeasibility on Iteration~1}

Constraint~\eqref{eq:eq-hard} requires all $n$ linearized shed values to be exactly equal.
This imposes $n-1$ independent linear equations on the $n$ weight perturbations $\Delta w_j$.
The system is infeasible when:
\begin{enumerate}
    \item The Jacobian $\mathbf{J}$ is rank-deficient (e.g., loads within the same load block have identical sensitivity rows), making it structurally impossible to independently equalize all shed values.
    \item The required weight perturbation lies outside the feasible region defined by~\eqref{eq:wbounds}--\eqref{eq:trust}, even with full rank.
\end{enumerate}
Unlike the proportional fairness case, increasing the trust-region radius~$\Delta$ cannot resolve structural rank deficiency.

\subsection{Reformulation: Relaxed Inequality with Quadratic Penalty}

We relax the hard equality~\eqref{eq:eq-hard} to an inequality and add a quadratic penalty that encourages equal shedding:
\begin{align}
    \min_{\mathbf{w},\,t} \quad & t + \sum_{i=1}^{n} \bigl(\hat{p}_i^{\mathrm{shed}}(\mathbf{w}) - t\bigr)^2 \label{eq:eq-obj}\\
    \text{s.t.} \quad
    & t \ge \hat{p}_i^{\mathrm{shed}}(\mathbf{w}), & i &= 1,\dots,n, \label{eq:eq-ineq}\\
    & t \ge 0, \notag \\
    & \eqref{eq:wbounds},\;\eqref{eq:trust}. \notag
\end{align}

The inequality~\eqref{eq:eq-ineq} ensures the problem is always feasible ($t$ can simply take the maximum shed value).
The quadratic penalty in~\eqref{eq:eq-obj} drives the shed values toward~$t$, recovering the equal-shedding behavior when the Jacobian permits it.
When exact equality is achievable, the penalty vanishes and the solution coincides with the original formulation.

\section{Summary of Changes}

\begin{table}[h]
\centering
\small
\begin{tabular}{@{}lll@{}}
\toprule
\textbf{Function} & \textbf{Issue} & \textbf{Fix} \\
\midrule
All functions & One-sided trust region $\Rightarrow$ degenerate solutions & Two-sided: $|w_j - w_j^{\mathrm{prev}}| \le \Delta$ \\
\texttt{proportional\_fairness} & Equality + bound infeasible; index mismatch & Shifted $\log(\cdot + c)$; aligned \texttt{pd} from caller \\
\texttt{min\_max\_load\_shed} & Non-smooth \texttt{maximum()} + tight $t\ge1$ & Epigraph: $t \ge \hat{p}_i^{\mathrm{shed}}\ \forall i$, $t \ge 0$ \\
\texttt{palma\_ratio} & Minimized ratio (wrong target); no efficiency & Target ratio $= 1$: $\min(\Sigma\mathrm{top} - \Sigma\mathrm{bot})^2$ \\
\texttt{equality\_min} & Hard equality structurally infeasible & Inequality + quadratic penalty \\
\bottomrule
\end{tabular}
\caption{Reformulations applied to upper-level fairness functions.}
\end{table}

\end{document}
